\documentclass[11pt]{article}
\usepackage{amsmath,amssymb,amsthm}
\usepackage[margin=1in]{geometry}
\newtheorem{theorem}{Theorem}
\newtheorem{lemma}{Lemma}
\begin{document}
\title{V108: An SCC-Cycle Bridge for Essential Ternary MUX Avoidance}
\author{NC0\_k-Avoid Laboratory}
\date{}
\maketitle

A signed essential MUX output has the form
\[
 g=o\oplus \operatorname{MUX}(s\oplus p_s,a\oplus p_a,b\oplus p_b).
\]
Fixing its output to a target exposes two exact binary clauses.  For either
branch one obtains an implication
\[
 x_s=\alpha\Longrightarrow x_d=\beta,
\]
where the source phase $\alpha$ is fixed by the selector polarity while the
arrival phase $\beta$ can be chosen by choosing the output target.

\begin{lemma}[Directed-cycle forcing]
Suppose selected branch arcs from pairwise distinct MUX outputs form a directed
simple cycle.  Targets can be chosen so that the induced implications transport
the source phase of each arc to the source phase of the next arc, except that
the final arc arrives at the complement of the initial phase.  Consequently
one literal on the initial selector is impossible and the opposite literal is
forced in every satisfying assignment.
\end{lemma}

\begin{theorem}[SCC-separated bridge]
Omit one MUX output $h$ and build the directed graph containing both branch arcs
of every remaining output.  Suppose the selector of $h$ lies on a directed
cycle in a cyclic strongly connected component $L$, and the selector literal
forced by that cycle activates a branch of $h$ whose data variable lies in a
different cyclic strongly connected component $R$.  Then a word outside the
range of the circuit can be constructed deterministically in polynomial time.
\end{theorem}

\begin{proof}
Target a cycle in $L$ to force the selector value that activates the chosen
branch of $h$.  Target a cycle through the arrival variable in $R$ so that this
variable is forced to the complement of its initial source phase.  Finally
target $h$ so its active branch forces the arrival variable to that initial
source phase.  The fixed-output 2-CNF is contradictory.  The two cycles use
disjoint outputs because every branch arc of one MUX has the same source
selector and $L,R$ are distinct SCCs.  SCC decomposition, cycle reconstruction,
and target assignment are polynomial-time operations.
\end{proof}

Let $\kappa_{\mathrm{SCC}}$ be the minimum number of outputs that must be ignored
before such a certificate appears.  Enumerating all ignored sets of size at most
$k$ gives runtime $O(m^k\operatorname{poly}(N))$; therefore constant
$\kappa_{\mathrm{SCC}}$ is polynomial-time avoidable.

\begin{theorem}[Strict separation]
For every $k\ge3$ there is an exact-stretch family with $n=2k$ inputs and
$m=2k+1$ essential canonical MUX outputs, consisting of two disjoint cyclic
families
\[
 \operatorname{MUX}(X_i,X_{i+1},X_{i+2})
\]
and one bridge $\operatorname{MUX}(L_0,L_2,R_0)$, such that
$\kappa_{\mathrm{SCC}}=0$.  The support family is Hall-minimal.  If $3\mid k$,
then its V102 strong-affine backdoor has exact size
\[
 \beta=\frac{4k}{3}=\frac{2n}{3}.
\]
\end{theorem}

\begin{proof}
After omitting the bridge, branch-zero arcs contain two disjoint directed cycles,
and the bridge joins the literal forced on the left to the right cycle.
Hall-minimality follows by an explicit shifted matching after deletion of any
cycle gate and the selector matching when the bridge is deleted.
For the backdoor bound, on each cycle a position outside the backdoor forces the
next two positions into it; hence at most $k/3$ positions are absent when
$3\mid k$.  A periodic pattern with one absent position per three attains this
bound while satisfying the bridge condition.
\end{proof}

\paragraph{Boundary.}
The theorem does not cover the residual MUX regime in which the useful cycles
remain in one cyclic SCC.  It does not solve unrestricted $NC^0_3$-Avoid or
P versus NP, and novelty or priority is not established.

\end{document}
